\documentclass[12pt]{article}
\usepackage{auck2e,multicol,amssymb,amsmath,color}
\usepackage{epsfig}
\usepackage[english]{babel}
\usepackage{blindtext}
\usepackage{microtype}
\usepackage{minted}

%\usepackage{graphicx}
\newcommand{\ds}{\displaystyle}
\newcommand{\mmark}[1]{[{#1 mark(s)}]}
\newcommand{\markk}[1]{[{\bf #1 marks}]}
\newcommand{\dspace}{\renewcommand{\baselinestretch}{2} \small \normalsize}
\newcommand{\sspace}{\renewcommand{\baselinestretch}{1} \small \normalsize}
\newcommand{\btheta}{\mbox{\boldmath $\theta$}}
\newcommand{\blambda}{\mbox{\boldmath $\lambda$}}
\newcommand{\bTheta}{\mbox{\boldmath $\Theta$}}
\newcommand{\bzeta}{\mbox{\boldmath $\zeta$}}
\newcommand{\bpi}{\mbox{\boldmath $\pi$}}
\newcommand{\bphi}{\mbox{\boldmath $\phi$}}
\newcommand{\ba}{\mbox{\boldmath $a$}}
\newcommand{\bb}{\mbox{\boldmath $b$}}
\newcommand{\bd}{\mbox{\boldmath $d$}}
\newcommand{\bm}{\mbox{\boldmath $m$}}
\newcommand{\bx}{\mbox{\boldmath $x$}}
\newcommand{\bh}{\mbox{\boldmath $h$}}
\newcommand{\bu}{\mbox{\boldmath $u$}}
\newcommand{\bP}{\mbox{\boldmath $P$}}
\newcommand{\bQ}{\mbox{\boldmath $Q$}}
\newcommand{\bR}{\mbox{\boldmath $R$}}
\newcommand{\by}{\mbox{\boldmath $y$}}
\newcommand{\bY}{\mbox{\boldmath $Y$}}
\newcommand{\bX}{\mbox{\boldmath $X$}}
\newcommand{\bU}{\mbox{\boldmath $U$}}
\newcommand{\bD}{\mbox{\boldmath $D$}}
\newcommand{\bV}{\mbox{\boldmath $V$}}
\newcommand{\bH}{\mbox{\boldmath $H$}}
\newcommand{\bgamma}{\mbox{\boldmath $\gamma$}}
\newcommand{\bGamma}{\mbox{\boldmath $\Gamma$}}
\newcommand{\balpha}{\mbox{\boldmath $\alpha$}}
\newcommand{\bbeta}{\mbox{\boldmath $\beta$}}
\newcommand{\bdelta}{\mbox{\boldmath $\delta$}}
\newcommand{\bomega}{\mbox{\boldmath $\omega$}}
\newcommand{\bpsi}{\mbox{\boldmath $\psi$}}
\newcommand{\bepsilon}{\mbox{\boldmath $\epsilon$}}
\newcommand{\be}{\begin{eqnarray}}
\newcommand{\ee}{\end{eqnarray}}
\def\BackgroundEPS#1#2#3#4{%
    \special{ps: @beginspecial @setspecial initmatrix
                0.9 setgray #2 #3 translate #4 dup scale}
    \special{ps: plotfile #1}
    \special{ps: @endspecial}
}
\newcommand{\wbugs}{{\sffamily\textbf{WinBUGS }}}
\newcommand{\coda}{{\sffamily\textbf{CODA }}}
\newcommand{\glm}{{\sffamily\textbf{G.L.M }}}
\newcommand{\rwb}{{\sffamily\textbf{R2WinBUGS }}}
\newcommand{\rar}{{\sffamily\textbf{Rock and roll man! }}}
\newcommand{\anova}{{\sffamily\textbf{ANOVA }}}
\newcommand{\mcmc}{{\sffamily\textbf{MCMC }}}
\newcommand{\dic}{{\sffamily\textbf{DIC }}}
\newcommand{\dagr}{{\sffamily\textbf{DAG }}}
\newcommand{\boxeq}[1]{%
\fboxrule=0.5mm%
\fboxsep=5mm%
\begin{center}\fcolorbox{black}{yellow}{\begin{minipage}{0.9\linewidth}
\begin{Beqnarray} #1 \end{Beqnarray}\end{minipage}}\end{center}%
\fboxrule=1pt%
\fboxsep=3pt%
}

\newcommand{\boxst}[1]{%
\fboxrule=0.5mm%
\fboxsep=5mm%
\begin{center}\fcolorbox{black}{yellow}{ \begin{minipage}{0.9\linewidth}#1 \end{minipage} }\end{center}%
\fboxrule=1pt%
\fboxsep=3pt%
}
%%%%%%%%%%%%%%%%%%%%%%%%%%%%%%%%%%%%%%%%%%%%%%%%%%%%%%%%%%%%%%%%%%%%%%%%
% Macro         :   \BackgroundEPS{filename}{tx}{ty}{s}
%-----------------------------------------------------------------------
% Description   :   Display an EPS file in the background of the
%                   current page.
%-----------------------------------------------------------------------
% Author        :   Bernard Desruisseaux (bernard@XYZZYcs.mcgill.ca)
%-----------------------------------------------------------------------
% Usage         :   \BackgroundEPS{pic.eps}{80}{100}{2}
%                   \BackgroundEPS{"`gunzip -c pic.eps.gz"}{80}{100}{2}
%                   \BackgroundEPS{"`zcat pic.eps.Z"}{80}{100}{2}
%%%%%%%%%%%%%%%%%%%%%%%%%%%%%%%%%%%%%%%%%%%%%%%%%%%%%%%%%%%%%%%%%%%%%%%%

\newcommand{\given}{\,|\,}

\paper{Introduction to Bayesian Statistics}
\due{Not for assessment}
\thetitle{\textbf{Tutorial 5}} 
\begin{document}


%\BackgroundEPS{bayes.eps}{0}{0}{1.2}
\makeband


\begin{Questions}
\item Remember the coal dataset and the changepoint model from lectures?
In this question you will do a posterior predictive check.

\begin{enumerate}
\item Download {\tt coal.R}, the data, and write up the changepoint model
from lectures.
\item Add extra code to the model for simulating entirely new data sets
for a posterior predictive check.
\item Plot at least 5 fake data sets from the posterior distribution. In your
opinion, does the real data look typical of the distribution of data simulated
from the model? i.e. does the model pass a posterior predictive check?
\end{enumerate}


\item Implement the coal changepoint model in R for Nested Sampling (this involves
editing the functions and variables in {\tt ns-model.R}. Use exactly the same
priors that we used before (i.e., log-uniform for $\lambda_1$ and $\lambda_2$,
uniform for the time of the changepoint.

\begin{enumerate}
\item [(a)] Run Nested Sampling. Inspect the diagnostic plots and the posterior
samples and make sure they look as you would expect. Also, make a note of the
marginal likelihood.
\item [(b)] The independent log-uniform priors for $\lambda_1$ and $\lambda_2$
are not really appropriate in part (a).
Realistically, we would expect that if there
is a change at some point, it will be a modest change (whereas under the independent
priors, you would most likely expect a drastic change).
Change the priors to
\begin{align}
\log(\lambda_1) &\sim \textnormal{Uniform}(-10, 10) \\
\log(\lambda_2) &\sim \textnormal{Normal}(\log(\lambda_1), 0.5^2).
\end{align}
Run this model. Do you get basically the same posterior distribution?
What about the marginal likelihood?
It should be higher than in
(a) --- why?
\item [(c)] Consider an alternative model where there is no changepoint,
so $\lambda_1$ applies to all years. Calculate the marginal likelihood for this
model too. Then work out the Bayes Factor for this model compared to the one
in (b)
\item [(d)] If the change-point model from (c) and the no-change model from
(b) have prior probabilities of 0.5 each, calculate their posterior
probabilities.
\end{enumerate}

\end{Questions}




\end{document}




